% Level 1, Week 4 -- workbook chapter
\chapterweek{4}{Keeping Score}{Variables are labelled boxes that remember things.}

\section{What you will learn}
\begin{itemize}[leftmargin=1.2em,itemsep=1pt]
  \item How a \textbf{variable} remembers a score or lives.
  \item The difference between \textbf{set} and \textbf{change}.
  \item How \textbf{pick random} makes your game full of surprises.
  \item How to build a catch game, and how to give \textbf{credit} for art you use.
\end{itemize}

\section{The big idea}
A \textbf{variable} is a box with a label. It holds \textbf{one} value at a time.

\begin{center}
\begin{tikzpicture}[>={Stealth},
  boxs/.style={draw,thick,minimum width=2.2cm,minimum height=1.3cm,fill=kcorange!15},
  lab/.style={font=\bfseries\small,fill=white,draw,inner sep=2pt}]
  \node[boxs] (s1) at (0,0) {\Large 0};
  \node[lab] at (s1.north) {score};
  \node[boxs] (s2) at (5.5,0) {\Large 1};
  \node[lab] at (s2.north) {score};
  \node[boxs] (s3) at (11,0) {\Large 2};
  \node[lab] at (s3.north) {score};
  \draw[->,thick] (s1.east) -- node[above,font=\small,align=center]{catch!\\change by 1} (s2.west);
  \draw[->,thick] (s2.east) -- node[above,font=\small,align=center]{catch!\\change by 1} (s3.west);
  \node[below=0.2cm of s1,font=\small,align=center] {start:\\set score to 0};
\end{tikzpicture}
\end{center}

\begin{itemize}[leftmargin=1.2em,itemsep=1pt]
  \item \textbf{set score to 0} empties the box and puts in 0. Use it once, at the start.
  \item \textbf{change score by 1} adds 1 to what is already in the box. Use it during the game.
  \item \textbf{pick random 1 to 6} is like rolling a die: you get a number, but you can't know which!
\end{itemize}

To make a variable: click \menu{Variables} $\rightarrow$ \menu{Make a Variable}, type a name, keep \enquote{For all sprites}, click \menu{OK}. A little score box appears on the stage.

\begin{exercise}[Pop the Balloon]
\begin{enumerate}[leftmargin=1.4em,itemsep=1pt]
  \item New project. Delete the cat. \menu{Choose a Sprite}: pick \textbf{Balloon1}.
  \item Make a variable called \texttt{pops}.
  \item Build both scripts. \enquote{pick random} is in \menu{Operators}: drop it into the x and y slots.
  \item Click the flag and pop the balloon! Then click the flag again: does \texttt{pops} start at 0?
\end{enumerate}
\begin{center}
\begin{scratch}[scale=0.8]
\blockinit{when \greenflag clicked}
\blockvariable{set \selectmenu{pops} to \ovalnum{0}}
\blockmove{go to x: \ovalnum{0} y: \ovalnum{0}}
\end{scratch}
\hspace{0.5cm}
\begin{scratch}[scale=0.8]
\blockinit{when this sprite clicked}
\blockvariable{change \selectmenu{pops} by \ovalnum{1}}
\blocksound{start sound \selectmenu{Pop}}
\blockmove{go to x: \ovaloperator{pick random \ovalnum{-200} to \ovalnum{200}} y: \ovaloperator{pick random \ovalnum{-140} to \ovalnum{140}}}
\end{scratch}
\end{center}
What happens if you take away \enquote{set pops to 0}? \answerlines{1}
\end{exercise}

\begin{exercise}[Apple Catch]
Catch the falling apples in a bowl!
\begin{enumerate}[leftmargin=1.4em,itemsep=1pt]
  \item New project. Delete the cat. Add the \textbf{Bowl} and the \textbf{Apple} sprites. Choose a backdrop you like.
  \item Make a variable called \texttt{score} (for all sprites).
  \item Click the \textbf{Bowl} and build its script. \enquote{mouse x} is in \menu{Sensing}.
  \item Click the \textbf{Apple} and build its script. \enquote{y position} is in \menu{Motion}; the \enquote{<} block is in \menu{Operators}.
  \item Play! Move the mouse left and right to catch apples.
\end{enumerate}
\begin{center}
\begin{tabular}{@{}c@{\hspace{1cm}}c@{}}
\begin{scratch}[scale=0.78]
\blockinit{when \greenflag clicked}
\blockmove{go to x: \ovalnum{0} y: \ovalnum{-140}}
\blockinfloop{forever}{
  \blockmove{set x to \ovalsensing{mouse x}}
}
\end{scratch}
&
\begin{scratch}[scale=0.78]
\blockinit{when \greenflag clicked}
\blockvariable{set \selectmenu{score} to \ovalnum{0}}
\blockmove{go to x: \ovaloperator{pick random \ovalnum{-200} to \ovalnum{200}} y: \ovalnum{170}}
\blockinfloop{forever}{
  \blockmove{change y by \ovalnum{-5}}
  \blockif{if \boolsensing{touching \ovalsensing*{Bowl} ?} then}{
    \blockvariable{change \selectmenu{score} by \ovalnum{1}}
    \blocksound{start sound \selectmenu{Chomp}}
    \blockmove{go to x: \ovaloperator{pick random \ovalnum{-200} to \ovalnum{200}} y: \ovalnum{170}}
  }
  \blockif{if \booloperator{\ovalmove{y position} < \ovalnum{-170}} then}{
    \blockmove{go to x: \ovaloperator{pick random \ovalnum{-200} to \ovalnum{200}} y: \ovalnum{170}}
  }
}
\end{scratch}
\\[4pt]
{\small\textbf{Bowl}} & {\small\textbf{Apple} (pick any sound in the dropdown)}
\end{tabular}
\end{center}
\textbf{Bug hunt!} Mei's apple script is the same, \textbf{but} it has no \enquote{go to x: (pick random) y: 170} inside the \enquote{touching Bowl?} block. Try it: take that block out. What happens to the score when you catch one apple? Why? \answerlines{2}
Put the block back. \textbf{Done when:} each apple you catch adds exactly 1, and missed apples start again at the top.
\end{exercise}

\begin{exercise}[Lives and Game Over (bonus)]
\begin{enumerate}[leftmargin=1.4em,itemsep=1pt]
  \item Make a variable \texttt{lives}. In the Apple's script, under \enquote{set score to 0}, add \enquote{set lives to 3}.
  \item Inside the \enquote{y position < $-170$} block (a missed apple), add \enquote{change lives by $-1$}.
  \item At the bottom of the forever loop, add this:
\end{enumerate}
\begin{center}
\begin{scratch}[scale=0.8]
\blockif{if \booloperator{\ovalvariable{lives} = \ovalnum{0}} then}{
  \blocklook{say \ovalnum{Game over!} for \ovalnum{2} seconds}
  \blockstop{stop \selectmenu{all}}
}
\end{scratch}
\end{center}
\textbf{Done when:} you start with 3 lives, lose one for each missed apple, and the game stops at 0.
\end{exercise}

\begin{safety}[Give credit]
The Apple, Bowl and Balloon were drawn by real people and shared in the Scratch library for everyone to use. Add a credits note: right-click the code area $\rightarrow$ \menu{Add Comment}, then type: \enquote{Credits: sprites and sounds from the Scratch library.} If you change a friend's project, that is a \textbf{remix}: say whose project you started from.
\end{safety}

\begin{vocab}
  \vocabword{variable}{A labelled box that remembers one value, like a score.}
  \vocabword{set}{Put a new value in the box (the old one is thrown away).}
  \vocabword{change}{Add to the value already in the box (use a minus number to take away).}
  \vocabword{random}{A surprise number you cannot predict, like rolling a die.}
  \vocabword{credit}{Saying who made the art, sounds or ideas you used.}
  \vocabword{remix}{Your own new version of someone else's project, with credit.}
\end{vocab}

\begin{reflection}
What is the difference between \enquote{set} and \enquote{change}?
\answerlines{2}
Whose art did you use today? How did you give credit?
\answerlines{1}
\end{reflection}
